% =====================================================================
% CHUONG 6 - KET LUAN VA HUONG PHAT TRIEN
% Tac gia: Bui Dinh Manh (Evaluation Engineer & Documentation)
% Cach dung: them % =====================================================================
% CHUONG 6 - KET LUAN VA HUONG PHAT TRIEN
% Tac gia: Bui Dinh Manh (Evaluation Engineer & Documentation)
% Cach dung: them % =====================================================================
% CHUONG 6 - KET LUAN VA HUONG PHAT TRIEN
% Tac gia: Bui Dinh Manh (Evaluation Engineer & Documentation)
% Cach dung: them % =====================================================================
% CHUONG 6 - KET LUAN VA HUONG PHAT TRIEN
% Tac gia: Bui Dinh Manh (Evaluation Engineer & Documentation)
% Cach dung: them \input{chuong6} vao main.tex.
% Neu dung documentclass{article}, doi \chapter -> \section.
% =====================================================================

\chapter{Kết luận và hướng phát triển}
\label{chap:conclusion}

\section{Kết luận}
Đề tài \textbf{"Đánh giá nguy cơ mắc bệnh tim bằng Mạng nơ-ron nhân tạo (ANN)"}
đã được thực hiện hoàn chỉnh theo đúng quy trình của một dự án học máy:
\begin{enumerate}
  \item \textbf{Dữ liệu:} thu thập bộ Heart Disease (Kaggle), tiền xử lý kỹ
        lưỡng -- phát hiện và loại bỏ \textbf{723 bản ghi trùng lặp} (còn 302
        bản duy nhất), xử lý ngoại lệ, chuẩn hoá đặc trưng và chia tập
        train/test 80/20.
  \item \textbf{Mô hình:} thiết kế và huấn luyện Mạng nơ-ron nhân tạo nhiều tầng
        ($13 \to 64 \to 32 \to 16 \to 1$) với ReLU, Sigmoid, Adam, Binary
        Crossentropy và Early Stopping.
  \item \textbf{Đánh giá:} mô hình đạt \textbf{Accuracy $\approx 0.74$},
        \textbf{F1-Score $\approx 0.76$} và \textbf{ROC-AUC $\approx 0.85$} trên
        tập test; thử nghiệm 4 cấu hình siêu tham số và chọn được mô hình tốt
        nhất (\textbf{ANN-3}, F1 $\approx 0.79$, Recall $\approx 0.82$).
\end{enumerate}

\textbf{Kết quả cho thấy:} mạng nơ-ron nhân tạo là một hướng tiếp cận
\textbf{khả thi và hiệu quả} cho bài toán sàng lọc nguy cơ bệnh tim từ dữ liệu
y tế dạng bảng, với khả năng phân tách hai lớp tốt (ROC-AUC $\approx 0.85$).

\section{Đóng góp của đề tài}
\begin{itemize}
  \item Xây dựng được một \textbf{pipeline hoàn chỉnh, tự động} từ dữ liệu thô
        đến kết quả đánh giá (\texttt{python main.py all}), có tổ chức mã nguồn
        rõ ràng, chú thích đầy đủ.
  \item Chỉ ra và xử lý đúng vấn đề \textbf{dữ liệu trùng lặp gây rò rỉ dữ liệu}
        -- một lỗi phổ biến dễ làm kết quả bị "ảo".
  \item Cung cấp bộ công cụ \textbf{đánh giá và trực quan hóa} đầy đủ (các chỉ
        số, Confusion Matrix, ROC, biểu đồ so sánh) phục vụ phân tích.
\end{itemize}

\section{Hạn chế}
\begin{itemize}
  \item \textbf{Kích thước dữ liệu nhỏ:} sau khi loại trùng chỉ còn 302 bản
        ghi, hạn chế khả năng học của mạng nơ-ron và làm kết quả dao động giữa
        các cấu hình.
  \item \textbf{Dữ liệu cũ (từ 1988)} và có một số giá trị bất thường về mã hoá
        (\texttt{ca = 4}, \texttt{thal = 0}), có thể không phản ánh đầy đủ thực
        tế hiện nay.
  \item Mô hình mới dừng ở mức \textbf{sàng lọc/hỗ trợ}, chưa thể thay thế chẩn
        đoán y khoa.
\end{itemize}

\section{Hướng phát triển}
\begin{itemize}
  \item \textbf{Mở rộng dữ liệu:} thu thập thêm dữ liệu thực tế, cập nhật và đa
        dạng hơn để tăng độ chính xác và khả năng tổng quát.
  \item \textbf{Thử nghiệm mô hình \& kỹ thuật khác:} so sánh ANN với các mô
        hình như Random Forest, XGBoost; áp dụng k-fold cross-validation, cân
        bằng dữ liệu, tinh chỉnh siêu tham số tự động (Grid/Random Search).
  \item \textbf{Giải thích mô hình (Explainable AI):} dùng SHAP/feature
        importance để giải thích yếu tố nào ảnh hưởng đến dự đoán -- quan trọng
        trong y tế.
  \item \textbf{Triển khai ứng dụng:} đóng gói mô hình thành web/app cho phép
        nhập chỉ số và trả về cảnh báo nguy cơ, hỗ trợ bác sĩ và người dùng.
\end{itemize}

\section{Lời kết}
Qua đề tài, nhóm đã vận dụng được kiến thức môn Trí tuệ nhân tạo vào một bài
toán y tế thực tiễn, nắm vững quy trình xây dựng -- huấn luyện -- đánh giá một
mô hình học máy, đồng thời rèn luyện kỹ năng làm việc nhóm và quản lý mã nguồn
bằng Git.
 vao main.tex.
% Neu dung documentclass{article}, doi \chapter -> \section.
% =====================================================================

\chapter{Kết luận và hướng phát triển}
\label{chap:conclusion}

\section{Kết luận}
Đề tài \textbf{"Đánh giá nguy cơ mắc bệnh tim bằng Mạng nơ-ron nhân tạo (ANN)"}
đã được thực hiện hoàn chỉnh theo đúng quy trình của một dự án học máy:
\begin{enumerate}
  \item \textbf{Dữ liệu:} thu thập bộ Heart Disease (Kaggle), tiền xử lý kỹ
        lưỡng -- phát hiện và loại bỏ \textbf{723 bản ghi trùng lặp} (còn 302
        bản duy nhất), xử lý ngoại lệ, chuẩn hoá đặc trưng và chia tập
        train/test 80/20.
  \item \textbf{Mô hình:} thiết kế và huấn luyện Mạng nơ-ron nhân tạo nhiều tầng
        ($13 \to 64 \to 32 \to 16 \to 1$) với ReLU, Sigmoid, Adam, Binary
        Crossentropy và Early Stopping.
  \item \textbf{Đánh giá:} mô hình đạt \textbf{Accuracy $\approx 0.74$},
        \textbf{F1-Score $\approx 0.76$} và \textbf{ROC-AUC $\approx 0.85$} trên
        tập test; thử nghiệm 4 cấu hình siêu tham số và chọn được mô hình tốt
        nhất (\textbf{ANN-3}, F1 $\approx 0.79$, Recall $\approx 0.82$).
\end{enumerate}

\textbf{Kết quả cho thấy:} mạng nơ-ron nhân tạo là một hướng tiếp cận
\textbf{khả thi và hiệu quả} cho bài toán sàng lọc nguy cơ bệnh tim từ dữ liệu
y tế dạng bảng, với khả năng phân tách hai lớp tốt (ROC-AUC $\approx 0.85$).

\section{Đóng góp của đề tài}
\begin{itemize}
  \item Xây dựng được một \textbf{pipeline hoàn chỉnh, tự động} từ dữ liệu thô
        đến kết quả đánh giá (\texttt{python main.py all}), có tổ chức mã nguồn
        rõ ràng, chú thích đầy đủ.
  \item Chỉ ra và xử lý đúng vấn đề \textbf{dữ liệu trùng lặp gây rò rỉ dữ liệu}
        -- một lỗi phổ biến dễ làm kết quả bị "ảo".
  \item Cung cấp bộ công cụ \textbf{đánh giá và trực quan hóa} đầy đủ (các chỉ
        số, Confusion Matrix, ROC, biểu đồ so sánh) phục vụ phân tích.
\end{itemize}

\section{Hạn chế}
\begin{itemize}
  \item \textbf{Kích thước dữ liệu nhỏ:} sau khi loại trùng chỉ còn 302 bản
        ghi, hạn chế khả năng học của mạng nơ-ron và làm kết quả dao động giữa
        các cấu hình.
  \item \textbf{Dữ liệu cũ (từ 1988)} và có một số giá trị bất thường về mã hoá
        (\texttt{ca = 4}, \texttt{thal = 0}), có thể không phản ánh đầy đủ thực
        tế hiện nay.
  \item Mô hình mới dừng ở mức \textbf{sàng lọc/hỗ trợ}, chưa thể thay thế chẩn
        đoán y khoa.
\end{itemize}

\section{Hướng phát triển}
\begin{itemize}
  \item \textbf{Mở rộng dữ liệu:} thu thập thêm dữ liệu thực tế, cập nhật và đa
        dạng hơn để tăng độ chính xác và khả năng tổng quát.
  \item \textbf{Thử nghiệm mô hình \& kỹ thuật khác:} so sánh ANN với các mô
        hình như Random Forest, XGBoost; áp dụng k-fold cross-validation, cân
        bằng dữ liệu, tinh chỉnh siêu tham số tự động (Grid/Random Search).
  \item \textbf{Giải thích mô hình (Explainable AI):} dùng SHAP/feature
        importance để giải thích yếu tố nào ảnh hưởng đến dự đoán -- quan trọng
        trong y tế.
  \item \textbf{Triển khai ứng dụng:} đóng gói mô hình thành web/app cho phép
        nhập chỉ số và trả về cảnh báo nguy cơ, hỗ trợ bác sĩ và người dùng.
\end{itemize}

\section{Lời kết}
Qua đề tài, nhóm đã vận dụng được kiến thức môn Trí tuệ nhân tạo vào một bài
toán y tế thực tiễn, nắm vững quy trình xây dựng -- huấn luyện -- đánh giá một
mô hình học máy, đồng thời rèn luyện kỹ năng làm việc nhóm và quản lý mã nguồn
bằng Git.
 vao main.tex.
% Neu dung documentclass{article}, doi \chapter -> \section.
% =====================================================================

\chapter{Kết luận và hướng phát triển}
\label{chap:conclusion}

\section{Kết luận}
Đề tài \textbf{"Đánh giá nguy cơ mắc bệnh tim bằng Mạng nơ-ron nhân tạo (ANN)"}
đã được thực hiện hoàn chỉnh theo đúng quy trình của một dự án học máy:
\begin{enumerate}
  \item \textbf{Dữ liệu:} thu thập bộ Heart Disease (Kaggle), tiền xử lý kỹ
        lưỡng -- phát hiện và loại bỏ \textbf{723 bản ghi trùng lặp} (còn 302
        bản duy nhất), xử lý ngoại lệ, chuẩn hoá đặc trưng và chia tập
        train/test 80/20.
  \item \textbf{Mô hình:} thiết kế và huấn luyện Mạng nơ-ron nhân tạo nhiều tầng
        ($13 \to 64 \to 32 \to 16 \to 1$) với ReLU, Sigmoid, Adam, Binary
        Crossentropy và Early Stopping.
  \item \textbf{Đánh giá:} mô hình đạt \textbf{Accuracy $\approx 0.74$},
        \textbf{F1-Score $\approx 0.76$} và \textbf{ROC-AUC $\approx 0.85$} trên
        tập test; thử nghiệm 4 cấu hình siêu tham số và chọn được mô hình tốt
        nhất (\textbf{ANN-3}, F1 $\approx 0.79$, Recall $\approx 0.82$).
\end{enumerate}

\textbf{Kết quả cho thấy:} mạng nơ-ron nhân tạo là một hướng tiếp cận
\textbf{khả thi và hiệu quả} cho bài toán sàng lọc nguy cơ bệnh tim từ dữ liệu
y tế dạng bảng, với khả năng phân tách hai lớp tốt (ROC-AUC $\approx 0.85$).

\section{Đóng góp của đề tài}
\begin{itemize}
  \item Xây dựng được một \textbf{pipeline hoàn chỉnh, tự động} từ dữ liệu thô
        đến kết quả đánh giá (\texttt{python main.py all}), có tổ chức mã nguồn
        rõ ràng, chú thích đầy đủ.
  \item Chỉ ra và xử lý đúng vấn đề \textbf{dữ liệu trùng lặp gây rò rỉ dữ liệu}
        -- một lỗi phổ biến dễ làm kết quả bị "ảo".
  \item Cung cấp bộ công cụ \textbf{đánh giá và trực quan hóa} đầy đủ (các chỉ
        số, Confusion Matrix, ROC, biểu đồ so sánh) phục vụ phân tích.
\end{itemize}

\section{Hạn chế}
\begin{itemize}
  \item \textbf{Kích thước dữ liệu nhỏ:} sau khi loại trùng chỉ còn 302 bản
        ghi, hạn chế khả năng học của mạng nơ-ron và làm kết quả dao động giữa
        các cấu hình.
  \item \textbf{Dữ liệu cũ (từ 1988)} và có một số giá trị bất thường về mã hoá
        (\texttt{ca = 4}, \texttt{thal = 0}), có thể không phản ánh đầy đủ thực
        tế hiện nay.
  \item Mô hình mới dừng ở mức \textbf{sàng lọc/hỗ trợ}, chưa thể thay thế chẩn
        đoán y khoa.
\end{itemize}

\section{Hướng phát triển}
\begin{itemize}
  \item \textbf{Mở rộng dữ liệu:} thu thập thêm dữ liệu thực tế, cập nhật và đa
        dạng hơn để tăng độ chính xác và khả năng tổng quát.
  \item \textbf{Thử nghiệm mô hình \& kỹ thuật khác:} so sánh ANN với các mô
        hình như Random Forest, XGBoost; áp dụng k-fold cross-validation, cân
        bằng dữ liệu, tinh chỉnh siêu tham số tự động (Grid/Random Search).
  \item \textbf{Giải thích mô hình (Explainable AI):} dùng SHAP/feature
        importance để giải thích yếu tố nào ảnh hưởng đến dự đoán -- quan trọng
        trong y tế.
  \item \textbf{Triển khai ứng dụng:} đóng gói mô hình thành web/app cho phép
        nhập chỉ số và trả về cảnh báo nguy cơ, hỗ trợ bác sĩ và người dùng.
\end{itemize}

\section{Lời kết}
Qua đề tài, nhóm đã vận dụng được kiến thức môn Trí tuệ nhân tạo vào một bài
toán y tế thực tiễn, nắm vững quy trình xây dựng -- huấn luyện -- đánh giá một
mô hình học máy, đồng thời rèn luyện kỹ năng làm việc nhóm và quản lý mã nguồn
bằng Git.
 vao main.tex.
% Neu dung documentclass{article}, doi \chapter -> \section.
% =====================================================================

\chapter{Kết luận và hướng phát triển}
\label{chap:conclusion}

\section{Kết luận}
Đề tài \textbf{"Đánh giá nguy cơ mắc bệnh tim bằng Mạng nơ-ron nhân tạo (ANN)"}
đã được thực hiện hoàn chỉnh theo đúng quy trình của một dự án học máy:
\begin{enumerate}
  \item \textbf{Dữ liệu:} thu thập bộ Heart Disease (Kaggle), tiền xử lý kỹ
        lưỡng -- phát hiện và loại bỏ \textbf{723 bản ghi trùng lặp} (còn 302
        bản duy nhất), xử lý ngoại lệ, chuẩn hoá đặc trưng và chia tập
        train/test 80/20.
  \item \textbf{Mô hình:} thiết kế và huấn luyện Mạng nơ-ron nhân tạo nhiều tầng
        ($13 \to 64 \to 32 \to 16 \to 1$) với ReLU, Sigmoid, Adam, Binary
        Crossentropy và Early Stopping.
  \item \textbf{Đánh giá:} mô hình đạt \textbf{Accuracy $\approx 0.74$},
        \textbf{F1-Score $\approx 0.76$} và \textbf{ROC-AUC $\approx 0.85$} trên
        tập test; thử nghiệm 4 cấu hình siêu tham số và chọn được mô hình tốt
        nhất (\textbf{ANN-3}, F1 $\approx 0.79$, Recall $\approx 0.82$).
\end{enumerate}

\textbf{Kết quả cho thấy:} mạng nơ-ron nhân tạo là một hướng tiếp cận
\textbf{khả thi và hiệu quả} cho bài toán sàng lọc nguy cơ bệnh tim từ dữ liệu
y tế dạng bảng, với khả năng phân tách hai lớp tốt (ROC-AUC $\approx 0.85$).

\section{Đóng góp của đề tài}
\begin{itemize}
  \item Xây dựng được một \textbf{pipeline hoàn chỉnh, tự động} từ dữ liệu thô
        đến kết quả đánh giá (\texttt{python main.py all}), có tổ chức mã nguồn
        rõ ràng, chú thích đầy đủ.
  \item Chỉ ra và xử lý đúng vấn đề \textbf{dữ liệu trùng lặp gây rò rỉ dữ liệu}
        -- một lỗi phổ biến dễ làm kết quả bị "ảo".
  \item Cung cấp bộ công cụ \textbf{đánh giá và trực quan hóa} đầy đủ (các chỉ
        số, Confusion Matrix, ROC, biểu đồ so sánh) phục vụ phân tích.
\end{itemize}

\section{Hạn chế}
\begin{itemize}
  \item \textbf{Kích thước dữ liệu nhỏ:} sau khi loại trùng chỉ còn 302 bản
        ghi, hạn chế khả năng học của mạng nơ-ron và làm kết quả dao động giữa
        các cấu hình.
  \item \textbf{Dữ liệu cũ (từ 1988)} và có một số giá trị bất thường về mã hoá
        (\texttt{ca = 4}, \texttt{thal = 0}), có thể không phản ánh đầy đủ thực
        tế hiện nay.
  \item Mô hình mới dừng ở mức \textbf{sàng lọc/hỗ trợ}, chưa thể thay thế chẩn
        đoán y khoa.
\end{itemize}

\section{Hướng phát triển}
\begin{itemize}
  \item \textbf{Mở rộng dữ liệu:} thu thập thêm dữ liệu thực tế, cập nhật và đa
        dạng hơn để tăng độ chính xác và khả năng tổng quát.
  \item \textbf{Thử nghiệm mô hình \& kỹ thuật khác:} so sánh ANN với các mô
        hình như Random Forest, XGBoost; áp dụng k-fold cross-validation, cân
        bằng dữ liệu, tinh chỉnh siêu tham số tự động (Grid/Random Search).
  \item \textbf{Giải thích mô hình (Explainable AI):} dùng SHAP/feature
        importance để giải thích yếu tố nào ảnh hưởng đến dự đoán -- quan trọng
        trong y tế.
  \item \textbf{Triển khai ứng dụng:} đóng gói mô hình thành web/app cho phép
        nhập chỉ số và trả về cảnh báo nguy cơ, hỗ trợ bác sĩ và người dùng.
\end{itemize}

\section{Lời kết}
Qua đề tài, nhóm đã vận dụng được kiến thức môn Trí tuệ nhân tạo vào một bài
toán y tế thực tiễn, nắm vững quy trình xây dựng -- huấn luyện -- đánh giá một
mô hình học máy, đồng thời rèn luyện kỹ năng làm việc nhóm và quản lý mã nguồn
bằng Git.
